\documentclass[12pt]{article}
\usepackage[utf8]{inputenc}
\usepackage[T1]{fontenc}
\usepackage{hyperref}
\title{Wolio Language Resources for K-12 Curriculum Development}
\author{Compiled by Claude Sonnet 5.5 (claude-sonnet-5-5) for Andr\'as Kornai}
\date{}
\begin{document}
\maketitle

\section{Introduction}
Wolio (ISO 639-3 wlo) is an Austronesian language of Baubau, Buton Island, Southeast Sulawesi. It was the court language of the Buton sultanate. Wikipedia gives about 65,000 speakers (2004).

\section{Primary Educational Resources}
No primers, textbooks or school programmes were found in search.

\section{Linguistic and Reference Materials}
\begin{itemize}
\item J.\ C.\ Anceaux, \emph{The Wolio Language} (grammar, 1952; second edition Foris/KITLV, 1988), and his Wolio dictionary (KITLV, 1987). Both are named in the Wikipedia article. I saw no direct URL for either.
\item A Wolio--Indonesian dictionary is available as a PDF from the Indonesian education ministry. It is linked from the Wikipedia article. I did not see the direct link.
\item Wikipedia, Wolio language: \url{https://en.wikipedia.org/wiki/Wolio_language}
\item Omniglot page on Wolio and its writing: \url{https://omniglot.com/writing/wolio.htm}
\end{itemize}

\section{Religious and Community Resources}
Nothing specific was found.

\section{Digital and Audio Resources}
The Endangered Languages Project has a Wolio page listing sources: \url{https://endangeredlanguages.com/node/112056}. I did not see any audio, video or software resources.

\section{Curriculum Development Context}
Anceaux's grammar and dictionary and the Indonesian-language Wolio dictionary are the main foundations. No school materials were found, and a curriculum would need new development.
\end{document}